\section{Arithmetic identities at the multiplicity boundary}
\label{sec:boundary-multiplicity}
This section expresses boundary multiplicities through the Euler
product, local conductor counts, and the derivative companion of
$\Xi$. The formulas identify the exact maps between arithmetic
residues, Fourier evaluations, global jets, and spectral trace defects.
They also locate the boundary estimates still needed for simplicity.

\subsection{A positive Poisson sum}
For $a>0$ and real $t$, define
\[
 P_a(t)=a\Re\frac{\xi'}{\xi}(\tfrac12+a+it).
\]
The centered even Hadamard product gives the absolutely convergent
identity
\begin{equation}\label{eq:poisson-multiplicity}
 P_a(t)=\sum_{\gamma\in\Gamma}
             m_\gamma\frac{a^2}{a^2+(t-\gamma)^2}.
\end{equation}
To verify its normalization, differentiate each paired factor
$\bigl(1+(s-1/2)^2/\gamma^2\bigr)^{m_\gamma}$.
Its logarithmic derivative is the sum of the fractions at
$1/2+i\gamma$ and $1/2-i\gamma$; their real parts give the two
terms displayed above. The zero count makes the tails locally
uniformly convergent for $a>0$.

At any zero ordinate $\gamma_0$ the corresponding term is exactly
$m_{\gamma_0}$. Every other term decreases as $a$ decreases to
zero, and its value at one fixed positive $a$ is a summable
majorant. Consequently
\begin{equation}\label{eq:poisson-infimum}
 \inf_{a>0}P_a(\gamma_0)
 =\lim_{a\downarrow0}P_a(\gamma_0)=m_{\gamma_0}.
\end{equation}
Thus simplicity at this ordinate is equivalent to the existence
of $a>0$ with $P_a(\gamma_0)<2$.

For $a>1/2$, put $s=1/2+a+it$. The absolutely convergent Euler
series gives exactly the same quantity as
\begin{equation}\label{eq:poisson-euler}
 P_a(t)=a\left[
 \Re\left\{\frac1s+\frac1{s-1}+\frac12\psi(s/2)\right\}
 -\frac12\log\pi
 -\sum_{n\ge2}\frac{\Lambda(n)}{n^{1/2+a}}\cos(t\log n)
 \right].
\end{equation}
The distance from the line matters quantitatively. The argument
bound $S(t)=O(\log t/\log\log t)$, available from the zero-location
theorem and \cite[Theorem~1, p.~2]{carneiro-chandee-milinovich2015},
implies
\begin{equation}\label{eq:poisson-exterior-size}
 P_{1/2}(t)=\tfrac14\log(t/(2\pi))
                +O(\log t/\log\log t).
\end{equation}
Here are the details of the passage to this formula. Integrate
the kernel $(1/4)/(1/4+(u-t)^2)$ against the positive zero-counting
measure. On $t/2\le u\le3t/2$, the main density is
$(2\pi)^{-1}\log(u/(2\pi))$; replacing its logarithm by
$\log(t/(2\pi))$ costs $O(\log t/t)$.
The integral of the kernel on the whole line is $\pi/2$.
Integration by parts against the argument error costs
$O(\log t/\log\log t)$, because the kernel has bounded total
variation. The remaining positive and negative tails cost
$O(\log t/t)$ by the zero count. This proves
\eqref{eq:poisson-exterior-size}.
Since \eqref{eq:poisson-multiplicity} increases with $a$, every
$a\ge1/2$ gives a value exceeding $2$ at all sufficiently large
heights. The upper threshold in \eqref{eq:poisson-infimum}
therefore requires control closer to the line. The logarithmic
derivative estimates of Chirre--Gon\c calves
\cite[Theorems~1--2 and (1.1)]{chirre-goncalves2022}
give explicit bounds in their stated distance range; no upper bound
below $2$ at every zero follows from the estimates used here.

\subsection{The prime ramp and its individual coefficients}
For $y>1$, define the arithmetic function
\[
 G(y)=\sum_{n\le y}\Lambda(n)(y/n-1)-y\log y+(1+\gamma_E)y.
\]
Direct integration of each summand, absolutely in $\Re s>1$,
gives
\begin{equation}\label{eq:prime-ramp-mellin}
 \int_1^\infty G(y)y^{-s-1}\,dy
 =-\frac{\zeta'(s)}{s(s-1)\zeta(s)}
       -\frac1{(s-1)^2}+\frac{1+\gamma_E}{s-1}.
\end{equation}
The pole at $s=1$ cancels, by
$\zeta'/\zeta=-1/(s-1)+\gamma_E+O(s-1)$.
At zero the residue is $\log(2\pi)$. At a nontrivial zero
$\rho$ it is $-m_\rho/[\rho(\rho-1)]$, and at $-2n$ it is
$-1/[2n(2n+1)]$. Mellin inversion therefore gives
\begin{equation}\label{eq:prime-ramp-zero-sum}
 G(y)=\log(2\pi)
  +\sum_{\gamma\in\Gamma}
       \frac{m_\gamma y^{1/2+i\gamma}}{1/4+\gamma^2}
  -\sum_{n\ge1}\frac{y^{-2n}}{2n(2n+1)}.
\end{equation}
For completeness, apply Mellin inversion first to
\[
 A(y)=\sum_{n\le y}\Lambda(n)(y/n-1),\qquad
 \int_1^\infty A(y)y^{-s-1}\,dy
       =-\frac{\zeta'(s)}{s(s-1)\zeta(s)}.
\]
On every fixed line $\Re s>1$, this transform is
$O((1+|\Im s|)^{-2})$, so its inversion is absolutely convergent.
At $s=1$ its Laurent terms are
$(s-1)^{-2}-(1+\gamma_E)(s-1)^{-1}$; after multiplication
by $y^s$, their residue is
$y\log y-(1+\gamma_E)y$.
Truncate the inversion at ordinates a distance
$\gg1/\log H$ from all zeros in the neighboring unit intervals.
The partial fractions and the unit-interval zero count give
$\zeta'/\zeta=O(\log^2H)$ on the horizontal portions of a fixed
strip. Their integrals tend to zero after division by $s(s-1)$.
Shift the left line through negative odd integers; the functional
equation bounds its logarithmic derivative by
$O(\log(|s|+2))$ away from the trivial zeros, and $y^{\Re s}$
makes its contribution tend to zero. The residue sums converge
absolutely: for the nontrivial zeros this follows from
$\sum m_\gamma/(1+\gamma^2)<\infty$, and for the trivial zeros
it follows from the displayed geometric factor. This justifies
the successive contour limits for $A(y)$. Subtracting the
explicit $s=1$ contribution, exactly as in the definition of
$G(y)$, proves \eqref{eq:prime-ramp-zero-sum}.

Put $R(u)=e^{-u/2}G(e^u)$. Its difference from
\[
 \sum_{\gamma\in\Gamma}
       \frac{m_\gamma}{1/4+\gamma^2}e^{i\gamma u}
\]
is integrable on $u\ge0$, by \eqref{eq:prime-ramp-zero-sum}.
The series converges uniformly and absolutely. Termwise averaging,
followed by dominated convergence over its summable coefficients,
therefore proves
\begin{equation}\label{eq:prime-ramp-coefficient}
 \lim_{U\to\infty}\frac1U\int_0^U
            R(u)e^{-i\gamma_0u}\,du
       =\frac{m_{\gamma_0}}{1/4+\gamma_0^2}.
\end{equation}
This is an exact map from the prime expression to each multiplicity.
An upper bound below $2/(1/4+\gamma_0^2)$ for this coefficient,
at every zero ordinate, would be precisely the simplicity
inequality. That arithmetic estimate remains open here.
The global identity
\[
 \sum_{\gamma\in\Gamma}\frac{m_\gamma}{1/4+\gamma^2}
       =2+\gamma_E-\log(4\pi)
\]
follows by evaluating twice the centered logarithmic derivative
at $s=1$. Its individual-coefficient consequence permits a bound
growing quadratically with $\gamma_0$ after the factor
$1/4+\gamma_0^2$ is restored.

One can see directly why a positive boundary residue and a
square-root size bound permit arbitrary integer multiplicities
in an analytic model. Fix $\gamma>0$, let $D=\gamma^2+1/4$,
and put
\[
 F_m(s)=\bigl(s(s-1)+D\bigr)^m,\qquad
 k_m=m/D,\qquad m=1,2,\ldots.
\]
Since $F_m'(1)/F_m(1)=m/D$, exact cancellation gives
\[
 -\frac{F_m'(s)/F_m(s)-sF_m'(1)/F_m(1)+k_m(1-s)}{s(s-1)}
 =\frac1D\frac{F_m'(s)}{F_m(s)}.
\]
The right side is the Laplace transform, initially for
$\Re s>1/2$, of
\[
 g_m(u)=\frac{2m}{D}e^{u/2}\cos(\gamma u).
\]
Its two residues are $m/D>0$ and
$|g_m(\log y)|\le(2m/D)\sqrt y$, while $F_m$ has order $m$
at both prescribed boundary zeros. These explicit formulas
verify the compatibility at every positive integer $m$.
The prime identity \eqref{eq:prime-ramp-mellin} supplies the
additional arithmetic structure in the zeta problem.

\subsection{Local zero counts for growing conductor}
\label{subsec:conductor-local-count}
The accepted Dirichlet zero-location input gives a local obstruction
to approximation across the critical line that also allows
holomorphic normalizing factors.

\begin{proposition}\label{prop:conductor-local-count}
Fix $t_0\in\mathbb R$ and $r>0$. For a primitive nonprincipal
character $\chi$ of conductor $q$, let $n_\chi(t_0,r)$ count the
nontrivial zeros, with their analytic multiplicities, whose
ordinates satisfy $|\gamma-t_0|<r$. There is a finite constant
$C_{t_0,r}$, independent of $q$ and $\chi$, such that
\[
 n_\chi(t_0,r)\ge \frac r{12}\log q-C_{t_0,r}.
\]
Let $U$ be an open neighborhood of $s_0=1/2+it_0$, let
$q_j\to\infty$ be conductors of primitive nonprincipal
characters $\chi_j$, and let $a_j$ be holomorphic and zero-free
on $U$. Every locally uniform holomorphic limit on $U$ of
$a_j(s)L(s,\chi_j)$ vanishes identically on the connected
component of $U$ containing $s_0$.
\end{proposition}

\begin{proof}
For the Fourier convention
$\widehat h(v)=\int_{\mathbb R}h(t)e^{-itv}\,dt$, put
\[
 w(y)=(1-y^2)\left(\frac{\sin y}{y}\right)^4,\qquad
 h(t)=w((t-t_0)/r).
\]
Removable values are understood. On the real line,
$h\le{\bf1}_{\{|t-t_0|<r\}}$: inside the interval its value
lies between zero and one, outside it is nonpositive, and at
the two endpoints it is zero. Moreover $\|h\|_1\le4r$, since
$|w(y)|\le1$ for $|y|\le1$ and $|w(y)|\le y^{-2}$ otherwise.
It is entire of exponential type $4/r$ and has decay
$O_{r,t_0}((1+|\Re z|)^{-2})$ on every fixed horizontal strip.

The transform identity
\[
 \widehat{(\sin y/y)^2}(v)=\frac\pi2(2-|v|)_+
\]
follows by convolving the transforms in
$\sin y/y=\frac12\int_{-1}^1e^{iyu}\,du$.
Plancherel, applied also to the weak derivative of this
triangle, gives
\[
 \int_{\mathbb R}\left(\frac{\sin y}{y}\right)^4dy
   =\frac{2\pi}{3},\qquad
 \int_{\mathbb R}y^2\left(\frac{\sin y}{y}\right)^4dy
   =\frac\pi2.
\]
Consequently
\[
 \int_{\mathbb R}h(t)\,dt=\frac{\pi r}{6},\qquad
 \operatorname{supp}\widehat h\subset[-4/r,4/r].
\]
The support assertion follows because the fourth power has
transform supported on $[-4,4]$, and multiplication by
$1-y^2$ adds its second transform derivative without
enlarging the support.

Write $a=0$ or $1$ according as $\chi(-1)=1$ or $-1$, and
$\psi=\Gamma'/\Gamma$. The primitive Guinand--Weil formula
in this convention is
\begin{align*}
 \sum_\rho h((\rho-\tfrac12)/i)
 ={}&\frac{\log(q/\pi)}{2\pi}\int_{\mathbb R}h(t)\,dt\\
 &+\frac1{2\pi}\int_{\mathbb R}h(t)
       \Re\psi\left(\frac14+\frac a2+\frac{it}{2}\right)dt\\
 &-\frac1{2\pi}\sum_{n\ge2}\frac{\Lambda(n)}{\sqrt n}
   \left\{\chi(n)\widehat h(\log n)
       +\overline{\chi(n)}\widehat h(-\log n)\right\}.
\end{align*}
See \cite[Section~2.1, equation~(2.1), p.~3]{chirre-pereira-delaat2021};
that source uses the Fourier exponential $e^{-2\pi ixt}$,
so its transform arguments $\pm\log n/(2\pi)$ become the
arguments displayed here. The conductor factor follows
from $(q/\pi)^{(s+a)/2}$ in the completion; reflection by
the functional equation gives the conjugate-character
term. There is no principal-character pole term.

Every zero lies on the critical line by the accepted input.
The zero sum is absolutely convergent by the decay of $h$
and the standard zero count. Its value is at most
$n_\chi(t_0,r)$. Only $n\le e^{4/r}$ occur in the prime sum,
whose absolute value is at most
\[
 \frac{4r}{\pi}\sum_{2\le n\le e^{4/r}}
                       \frac{\Lambda(n)}{\sqrt n}.
\]
The two possible gamma integrals are bounded independently
of $q$ and $\chi$, because
$\Re\psi(\alpha+it/2)=O(1+\log(2+|t|))$ for
$\alpha=1/4,3/4$. For example, this estimate follows from
\[
 \Re\psi(\alpha+iy)=\psi(\alpha)+
 \sum_{n\ge0}\frac{y^2}
 {(n+\alpha)((n+\alpha)^2+y^2)}
\]
by splitting at $n\asymp |y|$. Thus the formula gives
$\sum_\rho h(\gamma)=(r/12)\log q+O_{t_0,r}(1)$, proving
the stated bound. The constant is uniform in the conductor;
the radius is fixed throughout.

For the limit assertion, let $g$ be such a limit and suppose
its restriction to the component containing $s_0$ is not
identically zero. Choose $R>0$ with
$\overline{D(s_0,R)}\subset U$ and with no zero of $g$ on
the boundary. Isolatedness of the zeros permits this choice.
Local uniform convergence and Rouch\'e's theorem give a
fixed finite number of zeros in this disk for
$a_jL(\,\cdot\,,\chi_j)$ when $j$ is large. The factors
$a_j$ preserve every zero order. The first assertion with
$r=R/2$ instead gives at least
$(R/24)\log q_j-C_{t_0,R/2}$ zero occurrences on a critical
segment inside the disk, which tends to infinity. This
contradiction proves the component-wise conclusion.
\end{proof}

\subsection{The derivative companion and its common divisor}
\label{subsec:hb-boundary}
Put $F(z)=\Xi(z)$ and retain the distinct signed ordinates
$\Gamma$ and multiplicities $m_\gamma$. The paired canonical
product, with $\sum_\Gamma m_\gamma/\gamma^2<\infty$, gives
\[
 -\Im\frac{F'(z)}{F(z)}
 =(\Im z)\sum_{\gamma\in\Gamma}
          \frac{m_\gamma}{|z-\gamma|^2}>0
 \qquad(\Im z>0).
\]
Hence $E=F+iF'$ satisfies the upper-half-plane
Hermite--Biehler inequality
\[
 |E(z)|^2-|E^\#(z)|^2
  =4|F(z)|^2\left(-\Im\frac{F'(z)}{F(z)}\right)>0,
 \qquad E^\#(z)=\overline{E(\bar z)}=F-iF'.
\]
Reflection exchanges $E$ and $E^\#$ and preserves their
real boundary norm.

For the de Branges norm
$\|H\|_{\mathcal H(E)}^2=\int_{\mathbb R}|H(x)/E(x)|^2dx$,
where both $H/E$ and $H^\#/E$ belong to the upper
half-plane Hardy space, the kernel is
\begin{align}
 K_E(w,z)
 &=\frac{\overline{F(w)}F'(z)-\overline{F'(w)}F(z)}
                {\pi(\bar w-z)}\notag\\
 &=\frac1\pi\sum_{\gamma\in\Gamma}m_\gamma
       \frac{F(z)}{z-\gamma}
       \overline{\frac{F(w)}{w-\gamma}}.             \label{eq:hb-kernel}
\end{align}
The first equality follows by expanding the usual numerator
$\overline{E(w)}E(z)-\overline{E^\#(w)}E^\#(z)$ and dividing
by $2\pi i(\bar w-z)$. Subtracting the two logarithmic
derivatives gives the second. The series converges locally,
with removable values filled in. In particular
\begin{equation}\label{eq:hb-laguerre}
 \pi K_E(x,x)=F'(x)^2-F(x)F''(x).
\end{equation}
Its equality case at a multiple zero is also the $n=1$
case of \cite[Theorem~1.3 and Section~4, pp.~3, 10--12]{sra-xi2026}.

Define the convergent paired product
\[
 G(z)=\prod_{\gamma>0}(1-z^2/\gamma^2)^{m_\gamma-1},
 \qquad A=F/G,\quad B=F'/G,\quad \widetilde E=A+iB.
\]
At an order-$m$ zero, write
$F(z)=c_\gamma(z-\gamma)^m+O((z-\gamma)^{m+1})$.
Then $G$ has order $m-1$ there, $B(\gamma)=mA'(\gamma)\ne0$,
and
\[
 F'^2-FF''=mc_\gamma^2(z-\gamma)^{2m-2}
                       +O((z-\gamma)^{2m-1}).
\]
Thus $\widetilde E$ has no real zero, $E=G\widetilde E$, and
\[
 K_E(w,z)=G(z)\overline{G(w)}K_{\widetilde E}(w,z),
 \qquad K_{\widetilde E}(\gamma,\gamma)
                =mA'(\gamma)^2/\pi.
\]
The spaces satisfy the isometric identity
$\mathcal H(E)=G\mathcal H(\widetilde E)$.
Indeed, an entire $H$ with $H/E\in L^2(\mathbb R)$ must
vanish to at least the order of each real zero of $E$:
an integer deficit of one produces a nonintegrable
$|x-\gamma|^{-2}$ singularity. Therefore $H/G$ is entire,
and its two Hardy quotients and norm agree by cancellation.
Multiplication by $G$ proves the reverse inclusion.

The model-space inner function and its angular data are
\[
 \Theta=\frac{F-iF'}{F+iF'}
       =\frac{\widetilde E^\#}{\widetilde E},\qquad
 \Theta(\gamma)=-1,\qquad \Theta'(\gamma)=-\frac{2i}{m_\gamma}.
\]
Here $K(\Theta)=H^2(\mathbb C_+)\ominus\Theta H^2(\mathbb C_+)$,
with the boundary $L^2(dx)$ norm, and $K_\Theta$ denotes its
reproducing kernel.
The last two identities follow by substituting
$F/F'=(z-\gamma)/m_\gamma+O((z-\gamma)^2)$.
In the boundary $L^2(dx)$ normalization,
\[
 K_\Theta(\gamma,\gamma)=\frac1{\pi m_\gamma},
 \qquad
 \mu_\Theta(\{\gamma\})=\frac{2\pi}{|\Theta'(\gamma)|}
                       =\pi m_\gamma.
\]
These are the Clark masses for parameter $-1$ in
\cite[Section~2.4, equation~(2.8), and Proposition~4.1,
pp.~6--7, 10--11]{suzuki-weil2025}.

\subsection{Sampling, jets, and the boundary limit}
The normalized model-space basis from
\cite[Proposition~4.1]{suzuki-weil2025} is
\[
 e_\gamma(z)=i\sqrt{\frac{m_\gamma}{\pi}}\,
                  \frac{F(z)}{E(z)(z-\gamma)},\qquad
 e_\gamma(\eta)=\frac{\delta_{\gamma\eta}}{\sqrt{\pi m_\gamma}}.
\]
The source uses $\mathcal F_+h(\gamma)=\int h(x)e^{i\gamma x}dx$.
For the present convention $\widehat h=\mathcal F_-h$, the exact
conversion is $\mathcal F_+(h(-\,\cdot\,))(\gamma)=\widehat h(\gamma)$.
Equivalently $\mathcal F_+h(-\gamma)=\widehat h(\gamma)$;
these are alternative unitary reflections since
$m_{-\gamma}=m_\gamma$.
The weighted sampling isometry and this basis yield the unitary
\[
 \mathcal Uh
 =E\sum_{\gamma\in\Gamma}\sqrt{m_\gamma}\widehat h(\gamma)e_\gamma
 =\frac{i}{\sqrt\pi}\sum_{\gamma\in\Gamma}
                m_\gamma\widehat h(\gamma)\frac{F(z)}{z-\gamma}.
\]
Cauchy--Schwarz and $\sum_\Gamma m_\gamma/\gamma^2<\infty$
give locally uniform convergence of this series, with each removable
value filled in. Cauchy's formula therefore permits derivative
evaluation term by term. At an order-$m$ zero every nontarget
summand vanishes to order $m$, and hence
\begin{equation}\label{eq:hb-evaluation-jet}
 (\mathcal Uh)^{(m-1)}(\gamma)
       =\frac{i}{\sqrt\pi}F^{(m)}(\gamma)\widehat h(\gamma),
 \qquad
 (\mathcal Uh/G)(\gamma)
       =\frac{imA'(\gamma)}{\sqrt\pi}\widehat h(\gamma).
\end{equation}
This explicitly identifies the Fourier evaluations with
the corresponding global jets.

The multiplication extension in
\cite[Section~6, p.~18]{suzuki-weil2025} has the distinct
eigenvalues $\gamma$, each once. Its inverse, transported by
this unitary, is $C=\operatorname{diag}(1/\gamma)$. Therefore
the counted logarithmic trace and the distinct spectral
trace satisfy the exact identity
\[
 S_k-\operatorname{tr}C^{2k}
   =\sum_{\gamma\in\Gamma}\frac{m_\gamma-1}{\gamma^{2k}}
   =\sum_{\gamma\in\Gamma}
           \frac{2/|\Theta'(\gamma)|-1}{\gamma^{2k}}.
\]
This is also the limit of the finite defects $D_{N,k}$.

For $\omega>0$ put $E_\omega(z)=F(z+i\omega)$ and let
$K_\omega$ be its de Branges kernel. Direct differentiation
and the logarithmic derivative give
\[
 K_\omega(x,x)=\frac{\omega}{\pi}|F(x+i\omega)|^2
  \sum_{\eta\in\Gamma}\frac{m_\eta}{(x-\eta)^2+\omega^2},
 \qquad \omega^{-1}K_\omega(w,z)\longrightarrow K_E(w,z)
\]
locally uniformly. At an order-$m$ zero the exact leading
scale is
\begin{equation}\label{eq:hb-shift-scale}
 K_\omega(\gamma,\gamma)
   =\frac{m|c_\gamma|^2}{\pi}\omega^{2m-1}
                           (1+O_\gamma(\omega^2)).
\end{equation}
For the error term, the remaining zero sum is bounded near
$\gamma$, and reality of the Taylor coefficients gives
$|F(\gamma+i\omega)|^2
=|c_\gamma|^2\omega^{2m}(1+O_\gamma(\omega^2))$.

Strict interior positivity has a precise boundary limitation.
For $U(z)=|E(z)|^2-|E^\#(z)|^2$, the fact that $F'$ also
has only real zeros gives
\[
 \begin{gathered}
 \Delta U=4\bigl(|E'|^2-|(E^\#)'|^2\bigr)\ge0\quad(\Im z>0),\\
 U(\gamma+iy)=4m|c_\gamma|^2y^{2m-1}+O(y^{2m}).
 \end{gathered}
\]
The real-zero assertion for $F'$ follows by differentiating
the locally convergent paired polynomial products for $F$,
using Rolle's theorem and then Hurwitz's theorem.
The inward derivative vanishes at a multiple zero. The
subharmonic sign does not give the positive-superharmonic
boundary minimum inference. Explicitly, $F(z)=z^2$ gives
$U(x+iy)=8y(x^2+y^2)>0$ in the upper half-plane and
$\Delta U=64y>0$, with zero inward derivative at the origin.
The positive-kernel surgery already proved in this paper
gives the same cubic behavior in \eqref{eq:hb-shift-scale}
at its imposed double ordinate.

For comparison with the arithmetic canonical construction,
\cite[equations~(2.4)--(2.13), Theorems~2.1--2.4,
pp.~6--9]{suzuki-hamiltonians2021} uses
$E_{\omega,\nu}=E_\omega^\nu$ with $\nu\omega>1$ for zeta,
and obtains a positive determinant-one Hamiltonian for
each fixed permitted pair under the accepted zero-location
input. At real points the power has the exact kernel effect
\[
 K_{\omega,\nu}(x,x)
       =\nu|E_\omega(x)|^{2\nu-2}K_\omega(x,x).
\]
The displayed source results retain their fixed-parameter
quantifiers. A positive linear lower bound for
$K_\omega(\gamma,\gamma)$ as $\omega\downarrow0$ at every
zero ordinate remains unproved here; by
\eqref{eq:hb-shift-scale} it is the original multiplicity
margin.
